Grandes Modelos de Linguagem (LLMs) são utilizados na geração de conteúdo e no apoio a decisões. O alinhamento por preferências pode criar tensões entre produzir respostas apreciadas pelo usuário e preservar a correção factual. Nesse contexto, \textit{sycophancy} designa comportamentos de concordância indevida com crenças ou sugestões do usuário; em diálogos com múltiplos turnos, contestações repetidas podem influenciar as respostas do modelo. Este trabalho propõe investigar a capitulação diante de sugestões incorretas e a aceitação de correções válidas, considerando também a criticidade potencial das consequências de uma resposta errada. Para isso, pretende-se desenvolver, após a análise metodológica dos trabalhos-base, um protocolo reproduzível com perguntas verificáveis, condições de contestação e registros das interações. As medidas candidatas incluem taxa de capitulação, turno da primeira mudança, frequência de mudanças e aceitação de correções válidas, com populações elegíveis definidas separadamente. A pesquisa buscará comparar comportamentos entre modelos e cenários, distinguindo criticidade de dificuldade da questão. Nesta etapa, o protocolo experimental permanece em elaboração e não são apresentados resultados de testes com LLMs.
